% FAC-EN-U0080
% Source: fac_body.tex lines 3403--3447 (EOF).
% Authority: PDF physical pp. 82--83; printed pp. 277--278.
% Terminal unit: Paris signature and complete bibliography [1]--[19].

\medskip
\noindent\textsc{Paris}

\renewcommand{\refname}{Bibliography}
\phantomsection\label{fac:bibliography}
\begin{thebibliography}{99}
\addcontentsline{toc}{section}{Bibliography}

\bibitem[1]{fac-bib-1}\label{fac:bib:1}
N. \textsc{Bourbaki}. Théorie des Ensembles. Paris (Hermann).

\bibitem[2]{fac-bib-2}\label{fac:bib:2}
H. \textsc{Cartan}. \textit{Séminaire E.N.S.}, 1950--1951.

\bibitem[3]{fac-bib-3}\label{fac:bib:3}
H. \textsc{Cartan}. \textit{Séminaire E.N.S.}, 1951--1952.

\bibitem[4]{fac-bib-4}\label{fac:bib:4}
H. \textsc{Cartan}. \textit{Séminaire E.N.S.}, 1953--1954.

\bibitem[5]{fac-bib-5}\label{fac:bib:5}
H. \textsc{Cartan}. \textit{Variétés analytiques complexes et
cohomologie}. Colloque de Bruxelles, (1953), pp.~41--55.

\bibitem[6]{fac-bib-6}\label{fac:bib:6}
H. \textsc{Cartan} and S. \textsc{Eilenberg}. Homological Algebra.
Princeton Math. Ser., no.~19.

\bibitem[7]{fac-bib-7}\label{fac:bib:7}
S. \textsc{Eilenberg} and N. E. \textsc{Steenrod}. Foundations of
Algebraic Topology. Princeton Math. Ser., no.~15.

\bibitem[8]{fac-bib-8}\label{fac:bib:8}
W. V. D. \textsc{Hodge} and D. \textsc{Pedoe}. Methods of Algebraic
Geometry. Cambridge Univ. Press.

\bibitem[9]{fac-bib-9}\label{fac:bib:9}
W. \textsc{Krull}. Idealtheorie. Ergebnisse IV-3. Berlin (Springer).

\bibitem[10]{fac-bib-10}\label{fac:bib:10}
J. \textsc{Leray}. \textit{L'anneau spectral et l'anneau filtré
d'homologie d'un espace localement compact et d'une application
continue}. J. Math. Pures App., 29, (1950), pp.~1--139.

\bibitem[11]{fac-bib-11}\label{fac:bib:11}
G. \textsc{de Rham}. \textit{Sur la division de formes et de courants
par une forme linéaire}. Comment. Math. Helv., 28, (1954), pp.~346--352.

\bibitem[12]{fac-bib-12}\label{fac:bib:12}
P. \textsc{Samuel}. \textit{Commutative Algebra} (Notes by D. Hertzig).
Cornell Univ., 1953.

\bibitem[13]{fac-bib-13}\label{fac:bib:13}
P. \textsc{Samuel}. Algèbre locale. Mém. Sci. Math., CXXIII, Paris, 1953.

\bibitem[14]{fac-bib-14}\label{fac:bib:14}
J-P. \textsc{Serre}. \textit{Groupes d'homotopie et classes de groupes
abéliens}. Ann. of Math., 58, (1953), pp.~258--294.

\bibitem[15]{fac-bib-15}\label{fac:bib:15}
J-P. \textsc{Serre}. \textit{Un théorème de dualité}. Comment. Math.
Helv., 29, (1955), pp.~9--26.

\bibitem[16]{fac-bib-16}\label{fac:bib:16}
A. \textsc{Weil}. Foundations of Algebraic Geometry. Colloq. XXIX.

\bibitem[17]{fac-bib-17}\label{fac:bib:17}
A. \textsc{Weil}. \textit{Fibre-spaces in Algebraic Geometry} (Notes by
A. Wallace). Chicago Univ., 1952.

\bibitem[18]{fac-bib-18}\label{fac:bib:18}
O. \textsc{Zariski}. \textit{The concept of a simple point of an abstract
algebraic variety}. Trans. Amer. Math. Soc., 62, (1947), pp.~1--52.

\bibitem[19]{fac-bib-19}\label{fac:bib:19}
O. \textsc{Zariski}. \textit{Complete linear systems on normal varieties
and a generalization of a lemma of Enriques--Severi}. Ann. of Math., 55,
(1952), pp.~552--592.

\end{thebibliography}
